\section{Additional Nechako subdaily reference qualification}
The 2019--2024 author archive supersedes its 2019--2021 version. Before opening temperature sequences, published site dates were used to select the only river or creek candidate whose metadata could span all four UTC years 2020--2023; named lakes and a spillway were excluded. The frozen criterion requires exactly 96 unique quarter-hour timestamps per UTC day, finite values, and author flags P or A. The download contains 32 sites; only the selected candidate sequence was opened. No year meets this criterion, so no recovery model was scored. Each year also contains at least one nonfinite value independently of whether other flagged readings are admitted. This site is on the same Nechako River as an earlier BC exploratory case and is not an independent climatic confirmation. The source is Kokoszka et al. (2025), DOI 10.5281/zenodo.15053907, CC BY 4.0.
\begin{table}[p]\centering\small
\caption{Frozen P/A flag and 96-quarter-hour qualification of the selected Nechako site. Incomplete days include missing or nonfinite measurements, rejected author flags, duplicates or off-grid timestamps. Author flags are recommended classifications rather than proof of error-free measurements.}\label{tab:nechako-reference}
\begin{tabular}{lrrrrr}\toprule
Year & Days & Complete & Incomplete & Nonfinite records & Rejected flags\\\midrule
2020 & 366 & 341 & 25 & 1 & 1918\\
2021 & 365 & 211 & 154 & 6523 & 13982\\
2022 & 365 & 125 & 240 & 21021 & 21889\\
2023 & 365 & 127 & 238 & 17076 & 20129\\
\bottomrule\end{tabular}
\end{table}
